\documentclass[11pt]{article}

\usepackage[T1]{fontenc}
\usepackage{lmodern}
\usepackage{amsmath,amssymb,amsthm}
\usepackage{booktabs}
\usepackage[margin=1in]{geometry}
\usepackage[hidelinks]{hyperref}
\usepackage{microtype}

\hypersetup{
  pdftitle={Local rigidity of component-triangle Cayley deformations at order-seven permutation Euler points},
  pdfauthor={Xiaojian Zeng}
}

\newtheorem{theorem}{Theorem}
\newtheorem{proposition}[theorem]{Proposition}
\newtheorem{corollary}[theorem]{Corollary}
\newtheorem{remark}[theorem]{Remark}
\newcommand{\R}{\mathbb{R}}
\newcommand{\Q}{\mathbb{Q}}
\newcommand{\SO}{\operatorname{SO}}
\newcommand{\norm}[1]{\lVert #1\rVert}

\title{Local rigidity of component-triangle Cayley deformations\\
at order-seven permutation Euler points}
\author{Xiaojian Zeng\\
\small PBC School of Finance, Tsinghua University, Beijing 100190, China\\
\small \texttt{xiaojianzeng@foxmail.com}}
\date{}

\begin{document}
\maketitle

\begin{abstract}
For an orthogonal matrix $Q=(q_{ij})$ of order seven, write
$D(Q)=\sum_i q_{ii}^2$ and $A(Q)=\sum_i q_{i,8-i}^2$.
We study a three-parameter Cayley deformation through an unsigned
permutation matrix $P$ satisfying $D(P)=A(P)=1$ and having no eigenvalue
$-1$.  The three skew coordinates join, one at a time, each pair among the
singleton and the two three-cycle components of $P$.  In every such family,
$P$ is an isolated real solution of $D=A=1$.  The decisive quadratic
coefficients are established by an exhaustive exact-arithmetic census of
$112\cdot81=9{,}072$ families and a separately implemented rational-arithmetic
check.  This is a local obstruction for the specified families only; the
existence of a proper Euler magic matrix of order seven remains unresolved.
\end{abstract}

\section{Context and statement}

An Euler magic matrix of order $n$ is an integer matrix $M$ such that
$MM^{\mathsf T}=\gamma I_n$ for a nonzero integer $\gamma$, while the sums
of the entry-squares on the main diagonal and antidiagonal are both
$\gamma$.  It is \emph{proper} when all $n^2$ entry-squares are distinct.
This terminology and the equivalent rational normalization in odd order are
reviewed by M\"uller and Kominers~\cite{Muller2026,Kominers2026}.  Recent
work gives proper constructions in orders five and six
~\cite{Kominers2026,Mehat2026}, while M\"uller constructs improper
permutation examples in every order at least four and a proper example in
order eight~\cite{Muller2026}.

For $Q=(q_{ij})\in\SO(7)$ put
\[
 D(Q)=\sum_{i=1}^7q_{ii}^2,
 \qquad
 A(Q)=\sum_{i=1}^7q_{i,8-i}^2.
\]
Thus the normalized Euler locus is
\[
 V_7=\{Q\in\SO(7):D(Q)=A(Q)=1\}.
\]
Kominers notes that the odd-order reduction and mirror-rotation method apply
in order seven, but that a suitable seed is still needed
~\cite[Section~5]{Kominers2026}.  The Cayley transform gives a different
coordinate system on the orthogonal group and is standard in the study of
rational orthogonal matrices~\cite{LiebeckOsborne1991,Muller2026}.  We use it
to identify a limitation of a sparse deformation of the permutation-seed
route.

For distinct $a,b\in\{1,\ldots,7\}$ define the elementary skew matrix
\[
 F_{ab}=E_{ab}-E_{ba}\quad(a<b),
 \qquad F_{ab}=F_{ba}\quad(a>b).
\]
The convention merely fixes a sign for the coordinate on the unordered edge
$\{a,b\}$.

\begin{theorem}[Component-triangle local rigidity]\label{thm:main}
Let $P$ be an unsigned $7\times7$ permutation matrix having exactly one
entry $1$ on the main diagonal, exactly one entry $1$ on the antidiagonal,
and no eigenvalue $-1$.  Its cycle type is $1+3+3$.  Write $\{s\}$ for the
singleton cycle and $C_1,C_2$ for the two three-cycles.  Choose arbitrary
$u,v\in C_1$ and $w,t\in C_2$, and put
\[
 S_P=(I-P)(I+P)^{-1},\qquad
 E=xF_{su}+yF_{vw}+zF_{st},
\]
\[
 S(x,y,z)=S_P+E,\qquad
 Q(x,y,z)=(I-S)(I+S)^{-1}.
\]
Then $(0,0,0)$ is an isolated real common zero of
\[
 D(Q(x,y,z))-1=0,
 \qquad
 A(Q(x,y,z))-1=0.
\]
In particular, there is no nonconstant real-continuous solution germ in this
family through $P$.  Nor can a rational parametrization, defined at a real
parameter value mapping to $P$, parametrize non-anchor Euler points in this
family arbitrarily close to $P$.
\end{theorem}

The theorem is strictly local and family-specific.  It does not exclude
deformations involving additional Cayley coordinates, other seeds, or a
rational algebraic component on which $P$ is an isolated real point.  In
particular, it does not decide whether a proper order-seven Euler magic
matrix exists.

\section{Proof}

We first record the expansion that makes the finite part of the argument
small.  A permutation matrix has eigenvalue $-1$ exactly when its permutation
has an even cycle.  The hypothesis on the main diagonal says that there is
exactly one fixed point.  Every remaining cycle is therefore odd and has
length at least three; since six letters remain, the only possibility is
$3+3$.

The absence of the eigenvalue $-1$ makes $I+P$ invertible.  Orthogonality of
$P$ gives $S_P^{\mathsf T}=-S_P$, and direct substitution gives
\[
 Q(0,0,0)=P.
\]
Set
\[
 B=(I+S_P)^{-1}=\frac{I+P}{2}.
\]
For $E=xF_{su}+yF_{vw}+zF_{st}$, the identity
$(I+S_P+E)=(I+S_P)(I+BE)$ yields
\begin{align}
 Q(x,y,z)
 &=2(I+BE)^{-1}B-I \notag\\
 &=P-2BEB+2BEBEB+O(\norm{E}^3).\label{eq:taylor}
\end{align}
Write $L=-2BEB$ and $H=2BEBEB$ for the homogeneous terms of degrees one
and two.

The cycle decomposition $\{s\}\sqcup C_1\sqcup C_2$ makes $B$ block
diagonal.  Each active edge of $E$ joins two different blocks.  Hence $L$
vanishes at every position in the support of $P$.  The linear terms of both
$D(Q)-1$ and $A(Q)-1$ consequently vanish.  Their quadratic terms are
\begin{align}
 [D(Q)-1]_2
 &=\sum_{i=1}^7L_{ii}^2
   +2\sum_{i=1}^7P_{ii}H_{ii},\label{eq:d2}\\
 [A(Q)-1]_2
 &=\sum_{i=1}^7L_{i,8-i}^2
   +2\sum_{i=1}^7P_{i,8-i}H_{i,8-i}.\label{eq:a2}
\end{align}

Exact evaluation of \eqref{eq:d2}--\eqref{eq:a2} over all allowed choices
gives
\begin{equation}
 [D(Q)-1]_2=-2x^2-2z^2,\label{eq:decisive}
\end{equation}
and the coefficient $c=[y^2][A(Q)-1]_2$ belongs to
\[
 \left\{-\frac12,-\frac14,\frac14,\frac12,\frac34,1\right\};
\]
in particular, $c\ne0$.  Section~\ref{sec:verification} describes the
exhaustive exact calculation.

It remains to explain why these quadratic terms prove isolation, rather than
only the absence of a first-order tangent.  Suppose that distinct common real
zeros $r_k=(x_k,y_k,z_k)$ tend to the origin.  Passing to a subsequence, the
unit directions $r_k/\norm{r_k}$ tend to a unit vector
$r=(x_0,y_0,z_0)$.  Divide the first equation by $\norm{r_k}^2$ and take the
limit.  Equation~\eqref{eq:decisive} gives
$-2x_0^2-2z_0^2=0$, so $x_0=z_0=0$.  Applying the same argument to the
second equation leaves $c y_0^2=0$.  Since $c\ne0$, also $y_0=0$, contrary
to $\norm{r}=1$.  Thus the origin is isolated.  The two stated consequences
follow from continuity; a rational map is continuous wherever it is defined.
This proves Theorem~\ref{thm:main}.

\section{Exact finite verification}\label{sec:verification}

The coefficient census uses integer and rational arithmetic only.  It
enumerates all $7!$ permutations, retaining precisely those with one fixed
point, one antidiagonal hit, and cycle lengths $1,3,3$.  There are $112$.
After ordering the two three-cycles by their least elements, it enumerates
\[
 3\cdot(3\cdot3)\cdot3=81
\]
choices of $(\{s,u\},\{v,w\},\{s,t\})$ per permutation.  Swapping the two
three-cycles only interchanges the roles of $x$ and $z$, so this convention
loses no family.  For each of the resulting $9{,}072$ cases, it forms
$B,L,H$ and evaluates \eqref{eq:d2}--\eqref{eq:a2} exactly.

The primary checker uses SymPy rational matrices.  A separately written
checker uses only Python's standard-library \texttt{fractions.Fraction},
enumerates the anchors and edge choices afresh, and recomputes just the
coefficients needed above.  Both report no failed case.  The latter gives the
following complete distribution of $c=[y^2][A(Q)-1]_2$.

This emphasis on exact rational certificates and deterministic verification
also appears in nearby low-dimensional matrix algebra: Zeng gives an exact
rational Hadamard-factorization certificate together with a standard-library
verifier~\cite{Zeng2026Zenodo22242517}.  That work does not concern Euler
magic matrices and supplies no evidence for the local isolation theorem here.
\begin{center}
\begin{tabular}{@{}cr@{}}
\toprule
$c$ & number of families \\
\midrule
$-1/2$ & 1,728 \\
$-1/4$ & 3,456 \\
$ 1/4$ & 1,728 \\
$ 1/2$ & 1,512 \\
$ 3/4$ &   432 \\
$ 1$   &   216 \\
\midrule
Total & 9,072 \\
\bottomrule
\end{tabular}
\end{center}
The supplementary files \path{cayley_triangle_slices.py} and
\path{cayley_triangle_slices.json} contain the primary census; the files
\path{verify_cayley_triangle_tangents.py} and
\path{cayley_triangle_fraction_check.json} contain the separate check and
its output.  The reported run used Python 3.10.12 and SymPy 1.14.0.
The mathematical conclusion uses only the coefficient census, not the
additional bounded-grid experiments recorded by the primary script.

\section*{Computational and AI assistance disclosure}

AI systems assisted with exploration of the Cayley families, implementation
and review of the exact-arithmetic checks, literature inspection, and
manuscript drafting.  The decisive finite claim was checked by two separately
implemented programs, one of which uses only standard-library rational
arithmetic.  The proof and the exact certificates, rather than any numerical
fit or model assessment, support the theorem.  Responsibility for errors
remains with the author.

\bibliographystyle{plain}
\bibliography{references}

\end{document}
